\section{NuggetRecall}
\label{sec:nuggetrecall}

\judge{NuggetRecall} makes two kinds of LLM call. One call per topic writes an
answer key. One call per response then checks that response against the whole key.
The number of calls therefore grows with topics plus responses, not with responses
times facts. All of rag26 takes 9,838 judging calls. Its 39 responses with no text
score zero without a call.

\subsection{An answer key from the request alone}

For each topic, the LLM receives the request, which is the topic's title, problem
statement and background in a fixed order. We ask it for 8 to 20
self-contained facts that a complete, high-quality response should contain, with
the most important first. Each fact carries a tag. VITAL means a response that
leaves the fact out is clearly incomplete. OKAY means the fact is useful
supporting detail. Figure~\ref{fig:key} shows part of one key.
Appendix~\ref{app:prompts} gives both prompts in full.

The model never sees a response at this step. The AutoJudge framework enforces
this. Our workflow declares that the key does not depend on responses, so none are
passed in. Keeping responses out has three effects. A verbose system cannot shape
the key. The key can hold a fact that every system missed, which a key mined from
the responses cannot. And because the key comes from the topic alone, the data
policy lets us read it for any topic.

The parser accepts numbered lines that carry a VITAL or OKAY tag. It drops
duplicates and stops at 20 facts. A model that ignores the requested range
therefore cannot inflate the recall denominator. If a reply has fewer than two
usable facts, we retry with a fixed suffix appended to the prompt. If every attempt
fails, the topic falls back to a title-only key, a single fact made from its title. That never
happened with Qwen3-30B, our main judge model (Section~\ref{sec:setup}), on any pilot or test set. With gpt-oss-20b it happened on 5
of 122 ragtime25 topics and on 3 of 103 ragtime26 topics.

\begin{figure}
  \small
  \fbox{\parbox{0.94\columnwidth}{%
    \textbf{Request} (rag25-gen, topic 464, shortened). I want a thorough
    understanding of what makes up a community, including its definitions in
    various contexts like science [\ldots] I'm also interested in related terms
    like grassroots organizations, how communities set boundaries and
    priorities, and their roles in important areas such as preparedness and
    nation-building.

    \smallskip
    \textbf{Answer key} (5 of 20 facts).\\[2pt]
    \textsc{vital}\quad In ecology, a community refers to all the interacting
    populations of different species living in a particular area and sharing
    resources.\\[2pt]
    \textsc{vital}\quad Grassroots organizations are community-driven groups
    formed by local individuals to address specific social, environmental, or
    political issues from the bottom up [\ldots]\\[2pt]
    \textsc{vital}\quad Communities establish boundaries through shared identity,
    cultural norms, geographic limits, or formal rules, which help define
    membership and exclusion.\\[2pt]
    \textsc{okay}\quad Community boundaries can be fluid and change over time due
    to migration, cultural shifts, or political reorganization.\\[2pt]
    \textsc{okay}\quad Preparedness efforts in communities often involve
    collaboration with government agencies, NGOs, and emergency services.}}
  \caption{Part of an answer key that Qwen3-30B wrote from the request alone. The
  full key has 8 VITAL and 12 OKAY facts. The topic is the first one in the
  rag25-gen topics file. We did not pick it by hand.}
  \label{fig:key}
\end{figure}

\subsection{Checking each response in one call}

A single call judges each response against the whole key. The prompt holds the
request, the response and the numbered facts. The response appears as its
sentences joined with spaces. It is capped at 40,000 characters, a length no pilot
response reached. Citations and cited documents are left out, so
\judge{NuggetRecall} does not judge attribution (whether the cited documents back
what the response says). The prompt asks for one label per fact. The label is
\emph{support} if the response clearly states the fact. It is \emph{partial} if the
response states part of it or hedges it. Anything else gets \emph{no}. The model is
told to judge meaning rather than wording and to use no outside knowledge. The
sampling temperature is 0.

A reply is rejected unless it labels at least 80\% of the facts, rounded down to a
whole number. We do not accept a shorter reply and fill its gaps with \emph{no}.
That would invent misses and quietly understate the system. In an accepted reply,
the few facts left unlabeled count as \emph{no}. A rejected reply is retried up to
twice with a fixed suffix. The suffix changes the prompt and with it the cache key,
the lookup key for a stored reply. Since the suffix is fixed, replay stays
deterministic. A response that is still unreadable after three attempts scores
zero. We log how many there are. A reasoning model writes out intermediate steps
before its final answer. So its reply may hold a visible reasoning block before
the labels. The parser strips that block.

\subsection{Measures}

Let $K_q$ be the facts in the key for topic $q$ and $V_q \subseteq K_q$ the vital
ones. Let $\ell_{r,k}$ be the label that response $r$ received for fact $k$.
Lenient credit is $c(\text{support})=1$, $c(\text{partial})=0.5$ and
$c(\text{no})=0$. Strict credit $c'$ gives 1 to support only. Then
\begin{align}
  \texttt{NR\_VITAL}(r,q) &= \frac{1}{|V_q|}\sum_{k\in V_q} c(\ell_{r,k}),\\
  \texttt{NR\_WEIGHTED}(r,q) &= \frac{\sum_{k\in K_q} w_k\, c(\ell_{r,k})}{\sum_{k\in K_q} w_k},
\end{align}
with $w_k=1$ for vital facts and $0.5$ for okay ones. \texttt{NR\_ALL} is
\texttt{NR\_VITAL} taken over all of $K_q$. \texttt{NR\_VITAL\_STRICT} is
\texttt{NR\_VITAL} with $c'$. If a key has no vital fact, vital recall falls back to
recall over all facts. A system's score is its mean over topics. A missing response
counts as zero.

The fifth measure, \texttt{NR\_RANK}, uses the same judgments but aggregates them
differently. On each topic, every run gets the share of the other runs it beats on
\texttt{NR\_VITAL}, with ties split evenly. This within-topic percentile is
\begin{equation}
  \label{eq:percentile}
  p_q(r) = \frac{|\{r' : s_q(r') < s_q(r)\}| + \tfrac{1}{2}\,|\{r'\neq r : s_q(r') = s_q(r)\}|}{R_q - 1},
\end{equation}
where $s_q$ is the score on topic $q$ and $R_q$ is the number of runs with a
response on that topic. \texttt{NR\_RANK} is the mean of $p_q(r)$ over topics. A
plain mean over topics lets an easy topic, where every system scores high, count as
much as one that separates systems. Ranking within each topic before averaging
removes that difficulty effect. This follows the spirit of score
standardization~\cite{webber2008standardization}, which turns each topic's scores
into $z$-scores using their mean and standard deviation. We use a percentile and
not a $z$-score because the framework fills a missing cell with 0. Zero is the
right floor for a percentile but the wrong one for a $z$-score.

\subsection{Room to think for reasoning models}
\label{sec:budget}

The organizers run submitted code with gpt-oss models, which reason before they
answer. Each call has a reply budget, the most tokens the model may generate. A
reasoning model that uses up its budget stops in mid-thought and never writes its
labels. Retrying with the same budget cannot help. So each retry doubles the
budget, from 4,096 tokens to 8,192 and then 16,384. That still fits a context of
32,768 tokens. The first attempt keeps its original budget, so existing cache
entries still match. Section~\ref{sec:failures} shows the effect.

\subsection{Caching and replay}

All calls go through the organizers' minima-llm client~\cite{minimallm}. Its SQLite
cache builds each key from the model, the messages, the temperature and the reply
budget. The endpoint is not part of the key. Responses are sorted by run id before
any prompt is built, so prompts are byte-identical from one execution of the judge
to the next. That is what lets TIRA replay the judge from its cache with the
endpoint switched off.
